\section{Предметная область}
\label{sec:Chapter1} \index{Chapter1}
\subsection{Сетевые процессорные устройства}
Сетевые процессорные устройства (СПУ) — это специализированные системы, предназначенные для обработки сетевого трафика на высоких скоростях.
Проще говоря, такие устройства должны очень быстро принимать пакеты,
анализировать их и принимать решение о дальнейшей обработке. В отличие
от традиционных сетевых устройств, где логика работы часто жёстко зафиксирована на уровне аппаратуры, программируемые СПУ позволяют изменять
алгоритмы обработки пакетов и адаптировать их под разные задачи. Отличительной чертой современных СПУ является наличие программируемой логики обработки пакетов, что позволяет адаптировать устройство под новые протоколы и сетевые функции без замены аппаратного обеспечения.

\subsection{Архитектура RMT}
Одной из архитектур эталонной для построения высокопроизводительных программируемых сетевых устройств является RMT
(Reconfigurable Match Tables). Основная идея этой архитектуры заключается
в использовании набора настраиваемых таблиц сопоставления (match tables),
которые применяются на разных стадиях конвейера обработки пакетов.
На каждой стадии на основе выбранного ключа выполняется поиск в одной или нескольких таблицах, где правила сопоставления определяют, какие действия должны быть выполнены над пакетом.
Эти действия могут включать изменение заголовков, выбор следующей стадии обработки или передачу пакета на выходной интерфейс. Благодаря такой
структуре архитектура остаётся гибкой и позволяет реализовывать различные алгоритмы обработки сетевого трафика.

Match-action стадии — последовательность одинаковых стадий. Каждая стадия содержит:
	\begin{enumerate}
		\item Таблицы точного поиска, реализованные на SRAM.
		\item Таблицы тернарного и префиксного поиска, реализованные на TCAM.
		\item Кроссбары для подачи произвольных полей из PHV на вход таблиц.
		\item Action-блоки для параллельного изменения полей PHV.
	\end{enumerate}

\subsection{Программа обработки пакетов и язык P4}
Для описания поведения программируемого коммутатора используется язык P4\allowbreak{} (Programming Protocol-independent Packet Processors). Программа на P4 определяет набор логических match-action таблиц, каждая из которых включает набор полей заголовка, по которым производится поиск и списком возможных действий для них. P4 является аппаратно-независимым языком: одна и та же программа может компилироваться для различных устройств.

\subsection{Граф зависимостей и виды зависимостей внутри таблиц P4}
При компиляции P4-программы в конвейер RMT строится граф зависимостей таблиц (Table Dependency Graph, TDG). Вершинами графа являются логические таблицы, а рёбрами — зависимости между таблицами. Корректное размещение таблиц по стадиям конвейера должно удовлетворять всем видам зависимостей,иначе обработка пакетов будет нарушена.
Выделяют четыре основных типа зависимостей:

\begin{description}
	\item[Match-зависимость] возникает, когда таблица \(A\) изменяет поле, а таблица \(B\) выполняет поиск по этому полю. Эта зависимость соответствует классической зависимости «чтение после записи». Для корректной работы необходимо, чтобы \(A\) полностью завершила выполнение своих действий до того, как \(B\) начнёт поиск. Поэтому, таблицы с match-зависимостью не могут располагаться в одной стадии конвейера.
	
	\item[Action-зависимость] появляется, когда две таблицы \(A\) и \(B\) изменяют одно и то же поле, и важен результат, полученный в таблице \(B\). При размещении допускается нахождение \(A\) и \(B\) в одной стадии.
	
	\item[Successor-зависимость] заключается в следующем: выполнение таблицы \(B\) зависит от результата таблицы \(A\). Таблицы с successor-зависимостью могут размещаться в одной стадии.
	
	\item[Обратная match-зависимость] возникает, когда таблица \(A\) выполняет поиск по полю, которое впоследствии изменяется таблицей \(B\). Также требует строгого разделения стадий: \(A\) должна быть размещена в стадии, строго меньшей, чем \(B\).
\end{description}
 